\documentclass[conference]{IEEEtran}
\IEEEoverridecommandlockouts
\usepackage{cite}
\usepackage{amsmath,amssymb,amsfonts}
\usepackage{algorithmic}
\usepackage{graphicx}
\usepackage{textcomp}
\usepackage{xcolor}
\usepackage{url}
\usepackage{booktabs}
\graphicspath{{./figures/}}

\begin{document}

\title{Sub-Millisecond Multi-Implementation Autoscaling for Edge Serverless Computing via Constrained Deep Reinforcement Learning}

\author{\IEEEauthorblockN{Anonymous Authors}}

\maketitle

\begin{abstract}
% Section 0: Bilingual / English Abstract to be finalized in drafting phase.
\end{abstract}

\begin{IEEEkeywords}
Serverless Computing, Edge Intelligence, Knative Autoscaling, Multi-Implementation Dispatching, Deep Reinforcement Learning, PPO.
\end{IEEEkeywords}

\section{Introduction}
% - Motivation: Edge serverless bottleneck heterogeneity (CPU, GPU, Uplink).
% - Limitation of single-implementation KPA and ScaleWave GA (156 ms decision latency, CPU burn).
% - Our core idea: Sub-millisecond PPO policy with physical action masking.
% - 3-bullet contribution statement.

\section{Background and Related Work}
% - Knative Serving, KPA, and Concurrency Target scaling.
% - Multi-Implementation Equivalence and the ScaleWave MMKP Formulation.
% - Deep Reinforcement Learning for Edge Resource Management.

\section{System Architecture and Problem Formulation}
% - Multi-Resource Micro-Edge Cluster Model.
% - Constrained Multi-Implementation Scaling Formulation.
% - Physical Action Masking and Safety Guarantees.
% - Deep RL Agent Architecture (Actor-Critic PPO with 17-dimensional state space).

\section{Experimental Methodology}
% - Multi-Dataset Real-World Traces (Shanghai Telecom, Azure Functions 2019, NASA Flash Crowd).
% - Testbed Setup and Physical Container Dynamics (Cgroups CPU Throttling, Network Saturation).
% - Baseline Algorithms: Offload-only, Local-only, Local-offload, ScaleWave GA.
% - Key Performance Metrics (Throughput, Tail Latency, Solve Overhead, Efficiency).

\section{Evaluation}
% - Section A: Physical Testbed Reproduction (PerCom 2026 Figure 7).
% - Section B: Request Throughput and Success Rates under High Concurrency.
% - Section C: End-to-End and Tail Latency CDF Analysis.
% - Section D: Optimization Overhead and Control-Plane Decision Latency (866x Speedup).
% - Section E: Resource Allocation and Saturation Profiling.
% - Section F: Multi-Dataset Robustness (Azure Functions and NASA HTTP Access Logs).
% - Section G: Ablation Study: Elastic Capacity vs. Contended Hardware Dynamics.

\section{Discussion and Limitations}
% - Control-plane vs Data-plane Metric Decoupling.
% - Offline Training vs Online Heuristic Trade-offs.
% - Scalability to High Implementation Cardinality.

\section{Conclusion}
% - Summary of findings and future research directions.

\bibliographystyle{IEEEtran}
\bibliography{ref}

\end{document}
